%% main.tex — FIXTURE B of the underclaiming/buried-lede regression pair.
%%
%% Issue #1046 / #1048. This paper is deliberately BOLD AND LABELED: it
%% stakes a synthesis/research-program claim in the title, the abstract,
%% and the opening paragraph, and labels each contribution as
%% demonstrated / derived / synthesis / conjecture. Its method, numbers,
%% ablations, seeds, tables, threats section, and bibliography are
%% IDENTICAL to Fixture A's
%% (../../../underclaiming-buried-lede/...) — only the framing differs.
%%
%% Everything here is SYNTHETIC — the organization, the traces, the
%% measurements, and every bibliography entry are invented for the
%% fixture. Do not cite any of it.
%%
%% The Method-through-Experiments span is byte-identical to Fixture A's
%% and is pinned by tests/test_paper_underclaiming_fixtures.py, which is
%% what makes the score divergence attributable to framing alone. Edit
%% both fixtures or neither.

\documentclass{anvil-paper}

\title{Build Latency Is a History Problem: Repository Churn as a First-Class Scheduling Signal}
\author{Author Names Withheld}
\date{\today}

\begin{document}
\maketitle

\begin{abstract}
We argue that a distributed build system should treat version-control history
as a scheduling input, not merely as metadata for code analysis. The argument
rests on a measured asymmetry: in a fourteen-week trace of 2.1 million build
invocations across three production repositories, 8.3\% of source files
account for 71.4\% of remote-cache misses, and membership in that minority is
predictable a day ahead from churn and co-change centrality with AUC 0.87. We
build a scheduler that ranks files by that prediction and speculatively warms
their downstream targets; across an eight-week deployment to 1{,}412
engineers it cuts median build latency 22.4\% and p90 latency 31.1\% for a
5.8\% cluster-CPU cost. Neither history mining nor speculative warming is new;
the claim is that their composition makes history a scheduling signal, and the
ablation supports it — churn alone yields 9.1\%, centrality alone 7.4\%,
together 22.4\%. We conjecture, and do not here demonstrate, that the same
signal governs test selection and runner placement.
\end{abstract}

\section{Introduction}
\label{sec:intro}

A distributed build system decides what to compute and where, and it makes
those decisions from the dependency graph alone. We argue that this is the
wrong input set: the dependency graph says what \emph{could} change, while
version-control history says what \emph{will}, and the second signal is both
available and predictive a day in advance. Build latency, on this reading, is
less a caching problem than a history problem, and repository history belongs
in the scheduler beside the dependency graph rather than in the offline
analysis tools where it currently lives.

The measured asymmetry that motivates the claim is sharp. Across a
fourteen-week trace of 2.1 million build invocations, 8.3\% of source files
account for 71.4\% of remote-cache misses (Table~\ref{tab:concentration}), and
next-day membership in that minority is predictable from history with AUC 0.87
against 0.71 for the recency baseline in use
today~\cite{lindqvist2021prefetch}. A scheduler that acts on that prediction
cuts median build latency 22.4\% at a 5.8\% CPU cost.

We state the contributions at the scope each one is actually supported at,
labeling which are demonstrated, which are derived, which are synthesis, and
which are conjecture:

\begin{itemize}
  \item \textbf{Demonstrated.} Misses concentrate: 8.3\% of files account for
        71.4\% of misses across 2.1 million invocations, and the effect holds
        in all three repositories (Section~\ref{sec:experiments}). A
        history-derived predictor attains AUC 0.87 on next-day membership, and
        a scheduler acting on it cuts median build latency 22.4\% and p90
        31.1\% across 612{,}000 builds in an eight-week randomized deployment.
  \item \textbf{Derived.} The ranking rule follows from the concentration
        measurement and the standard speculative-work cost model rather than
        from tuning: rank by predicted miss probability, admit until the CPU
        budget binds (Section~\ref{sec:method}).
  \item \textbf{Synthesis.} Co-change mining~\cite{ferreira2016cochange,
        imani2017graphcentrality} and speculative build
        warming~\cite{lindqvist2021prefetch, salcedo2023speculative} are both
        mature; neither ingredient is ours. What is ours is the composition —
        history as a live scheduler input rather than an offline report — and
        the evidence that the composition is not redundant: the two history
        signals select different files and their combination beats the sum of
        the parts (Table~\ref{tab:ablation}).
  \item \textbf{Conjecture.} We expect the same history signal to govern
        regression-test selection and CI-runner placement, because both are
        scheduling decisions over the same change distribution. We do not test
        this here. The falsifiable form: a history-ranked test-selection
        policy should beat a failure-history policy~\cite{bhattacharya2022testsel}
        on tail CI latency by a margin comparable to the 16-point AUC gap we
        report for cache warming; a null result there would confine our claim
        to cache warming.
\end{itemize}

Section~\ref{sec:related} positions the work against the caching, prefetching,
and mining literatures, Section~\ref{sec:method} describes the trace and the
predictor, Section~\ref{sec:experiments} reports the measurements and the
ablation, and Section~\ref{sec:discussion} states the limits of the claim.

\section{Related Work}
\label{sec:related}

\paragraph{Build caching and remote execution.}
Content-addressed caching for build actions is mature~\cite{okonkwo2019remote},
and the questions it has generated are capacity, eviction, and transport --- all
of which treat the miss set as exogenous. Our claim is orthogonal to the
caching mechanism and complementary to it: we do not change what is cached, we
change which work the scheduler starts before it is asked for.

\paragraph{Prefetching in continuous integration.}
Lindqvist and Osei~\cite{lindqvist2021prefetch} warm CI artifacts by recency,
and Salcedo et al.~\cite{salcedo2023speculative} study admission policies for
speculative work on shared clusters. We inherit the warming mechanism from the
first and the budget-controller shape from the second. The disagreement is
about the ranking signal: recency is a one-feature proxy for history, and
Section~\ref{sec:experiments} shows it leaves most of the available signal
unused (AUC 0.71 against 0.87).

\paragraph{Co-change mining.}
Ferreira and Adeyemi~\cite{ferreira2016cochange} introduced the co-change graph
for impact analysis and Nakamura et al.~\cite{nakamura2020churn} established
the churn windows we adopt; the centrality measure is the standard eigenvector
formulation~\cite{imani2017graphcentrality}. We take all three unchanged. This
literature reads history offline, to advise humans about risk. The
contribution here is not a better mining technique but the observation that
its output is accurate enough, and fresh enough, to drive an online scheduling
decision --- a claim about where history belongs in the system, which the
mining literature does not make and which our ablation is what tests.

\paragraph{History-aware test selection.}
Bhattacharya and Lund~\cite{bhattacharya2022testsel} select regression tests
from historical failure data --- the closest prior instance of history driving an
online decision, and the reason our conjecture in Section~\ref{sec:discussion}
is a conjecture rather than a surprise. Their signal is failure history for a
test-selection decision; ours is change history for a build-scheduling
decision, and we do not claim their result transfers to our setting or the
reverse without the experiment we propose there.

\section{Method}
\label{sec:method}

\subsection{Trace collection}

Each build invocation is a directed acyclic graph of actions; the build tool
emits one record per action with the action's cache key, the resolved input
hashes, the outcome (\texttt{hit}, \texttt{miss}, or \texttt{local}), and the
wall-clock duration. We collected every such record for fourteen consecutive
weeks across three repositories: R1, a 41k-file monorepo; R2, a 9.8k-file
service repository; and R3, a 6.2k-file embedded firmware repository. The
trace contains 2.1 million build invocations comprising 903 million action
records.

A miss is attributed to a source file when that file's content hash appears in
the action's resolved input set and differs from the hash present at the last
observed hit for the same action identity. Actions whose input set changed in
more than one source file are attributed fractionally, splitting the miss
evenly across the changed inputs; Section~\ref{sec:discussion} reports the
sensitivity of our results to that choice.

\subsection{Predictor}

For each source file and each day we compute four features: the count of
commits touching the file in the preceding 30 days; the eigenvector centrality
of the file in the 90-day co-change graph~\cite{imani2017graphcentrality};
the reverse-dependency fan-out of the file in the build graph; and the mean
diff size across the file's commits in the preceding 30 days. A logistic
regression predicts whether the file will appear in the next day's
top-decile miss set. We fit with L2 regularization ($\lambda = 0.01$), five-fold
cross-validation grouped by week to avoid temporal leakage, and a fixed seed
(20260214) recorded in the released artifact.

\subsection{Pre-warming}

The pre-warmer ranks files by predicted probability, takes the top $k$, and
speculatively builds their downstream targets during idle cluster capacity.
The parameter $k$ is set by a controller that holds speculative work within a
6\% cluster-CPU budget; the controller and its budget are described in the
released configuration. No developer-visible behavior changes: speculative
work is preemptible and is discarded when a real build arrives.

\section{Experiments}
\label{sec:experiments}

\subsection{Miss concentration}

Table~\ref{tab:concentration} reports the concentration of the miss set. Across
the pooled trace, 8.3\% of source files account for 71.4\% of remote-cache
misses. Concentration is present in all three repositories and is strongest in
the monorepo R1, where 6.1\% of files account for 74.8\% of misses.

\begin{table}[t]
  \centering
  \begin{tabular}{lrrr}
    \toprule
    Repository & Files & Miss-heavy files (\%) & Misses covered (\%) \\
    \midrule
    R1 (monorepo)  & 41{,}200 & 6.1 & 74.8 \\
    R2 (service)   & 9{,}800  & 9.7 & 68.2 \\
    R3 (firmware)  & 6{,}200  & 11.4 & 66.9 \\
    \midrule
    Pooled         & 57{,}200 & 8.3 & 71.4 \\
    \bottomrule
  \end{tabular}
  \caption{Concentration of remote-cache misses. ``Miss-heavy files'' is the
  smallest set of files, ranked by attributed misses, covering the reported
  share of all misses over the fourteen-week trace.}
  \label{tab:concentration}
\end{table}

\subsection{Predictability}

The logistic regression attains AUC 0.87 (95\% CI [0.86, 0.88], five-fold
grouped cross-validation) on next-day top-decile miss-set membership. The
recency-only baseline of Lindqvist and Osei~\cite{lindqvist2021prefetch}
attains AUC 0.71 on the same folds. Table~\ref{tab:ablation} reports the
feature ablation.

\begin{table}[t]
  \centering
  \begin{tabular}{lrr}
    \toprule
    Signal & AUC & Median latency reduction (\%) \\
    \midrule
    Recency only (baseline)      & 0.71 & 5.3 \\
    Churn only                   & 0.79 & 9.1 \\
    Centrality only              & 0.76 & 7.4 \\
    Churn $+$ centrality (full)  & 0.87 & 22.4 \\
    \bottomrule
  \end{tabular}
  \caption{Feature ablation. Latency reductions are medians over the
  deployment arm described in Section~\ref{sec:experiments}.}
  \label{tab:ablation}
\end{table}

\subsection{Deployment}

We deployed the pre-warmer for eight weeks, cluster-randomized by engineer
across 1{,}412 engineers and 612{,}000 builds. Median end-to-end build latency
fell 22.4\% (95\% CI [20.9, 23.8]) in the treatment arm; p90 latency fell
31.1\%. Cluster CPU rose 5.8\%, within the configured budget. No regression
was observed in build correctness auditing, which compares speculative outputs
against the subsequently computed real outputs on a 1\% sample.

The combined signal outperforms the sum of its parts: churn alone yields 9.1\%
and centrality alone 7.4\%, while the combination yields 22.4\%. The two
signals select different files. Churn ranks recently edited leaves highly;
centrality ranks stable files whose edits propagate widely. Neither ordering
alone covers the miss set that drives tail latency.

\section{Discussion}
\label{sec:discussion}

\paragraph{What the claim rests on, and where it stops.}
The demonstrated part of the claim is bounded by the deployment: one
organization, three repositories, one build tool, eight weeks. The scheduling
argument --- that history belongs beside the dependency graph --- is supported by
the non-redundant ablation but is not proven by it, and a build environment
whose change distribution is flat rather than concentrated would not reproduce
the effect. The extension to test selection and runner placement is conjecture
and is labeled as such wherever it appears.

\paragraph{Threats to validity.}
The three repositories belong to one organization and share a CI lineage, so
practices such as branch cadence are held constant rather than varied. The
deployment arm ran on the cluster that produced the observational trace, so
the two are not independent. Fractional miss attribution for multi-input
changes is a modeling choice; re-running the analysis with
first-changed-input attribution moves pooled coverage from 71.4\% to 68.9\%
and does not change the ordering in Table~\ref{tab:ablation}. Finally, the
deployment measured wall-clock latency under the cluster's ordinary contention
and not under saturation, where speculative work is likelier to displace real
work.

\paragraph{Cost.}
The 5.8\% CPU overhead is a real cost, and organizations whose clusters run
near saturation should expect the trade to be less favorable than it was here.

\section{Conclusion}
\label{sec:conclusion}

Version-control history predicts, a day in advance and with AUC 0.87, which
files will drive a build cluster's cache misses, and a scheduler that acts on
that prediction cuts median build latency 22.4\% for a 5.8\% CPU cost. We take
this as evidence for the broader claim that history is a scheduling signal
rather than an offline analysis product --- a claim whose next test is not
another cache experiment but the test-selection experiment named in
Section~\ref{sec:discussion}.

\bibliographystyle{plainnat}
\bibliography{refs}

\end{document}
